\documentclass[a4paper,12pt]{article}
\usepackage{graphicx}
\usepackage[a4paper, top=1cm, bottom=2cm, left=2cm, right=2cm]{geometry}
\usepackage{xcolor}
\usepackage{hyperref}
\usepackage{eso-pic}

\usepackage[french]{babel}
\usepackage[T1]{fontenc}
\usepackage{mathrsfs}
\usepackage{amsmath}
\usepackage{amsfonts}
\usepackage{amssymb}

\usepackage{tikz}
\usetikzlibrary{arrows, shapes.geometric, fit}

\newcounter{correction}

\newcommand{\question}{%
  \refstepcounter{correction}%
  \textbf{\textcolor{black}{Question \thecorrection :}} \ignorespaces
}

\AddToShipoutPicture{
    \AtTextCenter{%
        \makebox[0pt]{\rotatebox{80}{\textcolor[gray]{0.6}{\fontsize{10cm}{10cm}\selectfont PGB}}}
    }
}

\begin{document}
\hrule
\begin{center}
    \begin{tabular}{@{} p{5cm} p{5cm} p{5cm} @{}}
        Lycée Dindéfélo & \quad\quad\quad Test 02 & 04 Novembre 2026 \\
    \end{tabular}
    \\[-0.01cm]
    \hrule
\end{center}
\begin{center}
    \textbf{\Large Racine Carrée et Nombres Réels} \\[0.2cm]
    \textbf{\large Professeur : M. BA} \\[0.2cm]
    \textbf{Classe : 3 ème} \\[0.2cm]
    \textbf{\small Durée : 15 minutes} \\[0.2cm]
    \textbf{\small Note :\quad\quad /6}
\end{center}

\textbf{\small Nom de l'élève :} \underline{\hspace{8cm}} \\[0.5cm]

\question (2 points) Simplifier ou calculer les expressions suivantes :
\[
\begin{aligned}
A&=\sqrt{144}\\
&=\\
&=\\
&\\
\end{aligned}
\quad\quad\quad
\begin{aligned}
B&=(\sqrt{13})^2\\
&=\\
&=\\
&\\
\end{aligned}
\quad\quad\quad
\begin{aligned}
C&=\sqrt{(-9)^2}\\
&=\\
&=\\
&\\
\end{aligned}
\quad\quad\quad
\begin{aligned}
D&=\sqrt{3^2 \times 5^2}\\
&=\\
&=\\
&=\\
\end{aligned}
\]

\vspace{0.8cm}

\question (2 points) Donner la valeur exacte, la valeur approchée et le cadre demandé pour $M = \sqrt{7}$ (sachant que $\sqrt{7} \approx 2,6457$) :
\begin{enumerate}
    \item[a)] Valeur exacte de $M =$ \underline{\hspace{4cm}}
    \item[b)] Valeur approchée de $M$ au centième près par défaut : \underline{\hspace{3cm}}
    \item[c)] Valeur approchée de $M$ au centième près par excès : \underline{\hspace{3cm}}
    \item[d)] Encadrement de $M$ à $10^{-2}$ près : \underline{\hspace{4cm}} $< \sqrt{7} <$ \underline{\hspace{4cm}}
\end{enumerate}

\vspace{0.8cm}

\question (2 points) Appliquer les propriétés des racines carrées pour calculer ou simplifier sous la forme d'un entier ou d'un quotient simple :
\[
\begin{aligned}
E&=\sqrt{2} \times \sqrt{32}\\
&=\\
&=\\
&=\\
&\\
\end{aligned}
\quad\quad\quad\quad\quad
\begin{aligned}
F&=\frac{\sqrt{75}}{\sqrt{3}}\\
&=\\
&=\\
&=\\
&\\
\end{aligned}
\quad\quad\quad\quad\quad
\begin{aligned}
G&=\sqrt{\frac{49}{81}}\\
&=\\
&=\\
&\\
\end{aligned}
\]

\end{document}